\chapter{开会，开会，开会}\label{s:meetings}

大多数人真的不擅长开会：
他们不做议程、
不做纪要、
他们啰嗦个没完、或跑题到不相干的事上、
他们说些陈词滥调、或重复别人说过的话，
只为了显得自己说了点什么，
还会开小会
（这几乎保证了这场会是在浪费时间）。
知道如何高效地主持会议，
是任何想把事情做成的人的核心技能；
而知道如何参与别人主持的会议，同样重要
（尽管它受到的关注少得多——正如一位同事所说，
人人都提供领导力培训，
但没人提供「追随力」培训）。

让会议高效的最重要规则并不神秘，
但很少有人遵守：

\begin{description}

\item[先判断是否真的需要开会。]
  如果唯一的目的只是分享信息，
  那就发一封简短邮件。
  记住，
  你读得比任何人讲得都快：
  如果某人有一些事实要让团队其他人吸收，
  最礼貌的传达方式就是把它写下来。

\item[写一份议程。]
  如果没人愿意为这场会写一份要点的讨论清单，
  这个会大概就不必开了。

\item[在议程里写上时间。]
  如果你在议程里写上每一项要花的时间，
  议程还能帮你防止前面的议题偷走后面议题的时间。
  你和任何新群体合作时的第一次估计都会过于乐观，
  所以后续会议把时间往上调。
  不过，
  你不该因为第一次超时，
  就安排第二次、第三次会：
  相反，
  要设法弄清你为什么会超时，并解决根本问题。

\item[排优先级。]
  每场会都是一个微型项目，
  所以工作应该像对待其他项目那样排优先级：
  影响大但花时间少的先做，
  花时间多但影响小的就跳过。

\item[指定一个人负责让会议往前走。]
  应该有一项任务交给某个人：让各项议题守时、
  提醒那些查邮件或开小会的人、
  让说得太多的人说到点子上、
  并请没发言的人表达意见。
  这个人\emph{不该}包办所有发言；
  事实上，
  在一场主持得好的会里，负责人的发言会比大多数其他参与者更少。

\item[要求礼貌。]
  谁都不许粗鲁，
  谁都不许东拉西扯，
  如果有人跑题，
  主持人有权也有责任说：
  「我们把那个留到别处讨论。」

\item[不许打断。]
  参与者如果想下一个发言，
  应该举一根手指、或贴一张便利贴。
  如果正在说话的人没注意到，
  主持人应该注意到。

\item[不用电子设备，]
  除非出于无障碍原因必须用。
  坚持让所有人把手机、平板和笔记本电脑调到礼貌模式
  （也就是合上）。

\item[做纪要。]
  应该由主持人之外的某个人做要点记录，
  记下分享的最重要的信息、
  做出的每一个决定、
  以及分配给某人的每一项任务。

\item[记笔记。]
  别人说话时，
  参与者应该记下自己想问的问题、或想表达的观点。
  （轮到你发言时，你会发现这让你显得多么聪明。）

\item[早点结束。]
  如果你的会定在 10:00—11:00，
  你应该争取在 10:50 结束，
  好给人时间在去下一个地方的路上上个厕所。

\end{description}

会议一结束，
就把纪要发邮件给所有人，或贴到网上：

\begin{description}

\item[没到场的人能跟上进展。]
  一个网页或一封邮件，是比问队友「我错过了什么」高效得多的跟进方式。

\item[每个人都能核对实际说了什么、承诺了什么。]
  不止一次，
  我翻看自己参加过的会议的纪要时会想：
  「我说过那句话吗？」
  或者，「等等，我没承诺过那时候把它弄好啊！」
  不管有意无意，
  人们记住的东西常常不一样；
  把它写下来，给了团队成员一个纠正错误的机会，
  这能在日后省下很多痛苦。

\item[人们能在后续会议上被问责。]
  如果你事后不跟进，
  列问题和行动项就毫无意义。
  如果你在用某种 issue 跟踪系统，
  就在会后立刻为每个新问题或新任务建一个 issue，
  并更新那些继续往下带的，
  然后每次会议一开始，就过一遍这些 issue 的清单。

\end{description}

\cite{Brow2007,Broo2016,Roge2018}在主持会议方面有很多建议。
依我的经验，
花一小时学怎么当主持人，
会是你做过的最好的投资之一。

\begin{aside}{便利贴与打断宾果}
  有些人太习惯听自己的声音了，
  不管屋里还有多少人，
  他们都非得讲掉一半时间。
  为了应对这一点，
  在会议开始时给每个人三张便利贴。
  他们每发一次言，
  就得取下一张便利贴。
  等他们的贴纸用完了，
  在每个人都至少用过一张之前，就不许再说话；
  到那时候，每个人又都拿回自己所有的便利贴。
  这保证了没人发言的频率超过最安静那个人的三倍，
  也彻底改变了大多数群体的动态：
  那些因为总被打断而早已放弃被听见的人，
  突然有了贡献的空间；
  而发言过于频繁的人则很快意识到自己一直有多不公平\footnote{
    这种事落到我头上时，我确实意识到了{\ldots}
  }。

  另一种技巧是打断宾果。
  画一张网格，用参与者的名字标注行和列。
  每当有人打断别人，
  就在相应的格子里加一笔，
  并在会议进行到一半时花点时间分享结果。
  大多数情况下，
  你会看到是一两个人在包办所有打断，
  而且往往自己都没意识到。
  单是这一点，往往就足以让他们收敛。
  （注意，这个技巧是为了管理打断，
  而不是发言时长：
  让对某个主题知道更多的人在会上的关于它的发言更多，也许是合适的，
  但反复打断别人，永远都不合适。）
\end{aside}

\seclbl{玛莎规则}{s:meetings-marthas-rules}

世界各地的组织都按照
\hreffoot{https://en.wikipedia.org/wiki/Robert\%27s\_Rules\_of\_Order}{《罗伯特议事规则》}开会，
但它比大多数小项目所需要的正式得多。
一种叫「玛莎规则」的轻量替代方案，
对基于共识的决策可能好得多~\cite{Mina1986}：

\begin{enumerate}

\item
  每次会议之前，
  任何想提案的人都可以把提案分享给全组来发起它。
  提案必须至少在会议前 24 小时提交，才能在那次会上被审议，
  并且必须包括：
  \begin{itemize}
  \item 一句话摘要；
  \item 提案全文；
  \item 任何必要的背景信息；
  \item 支持与反对的理由；以及
  \item 可能的替代方案
  \end{itemize}
  提案最长应为两页。

\item
  如果半数或以上的有表决权成员到场，会议即达到法定人数。

\item
  某人发起提案后，
  就由他负责这件事。
  除非发起人或其代表在场，否则全组不得讨论或表决这个议题。
  发起人还负责把这一项呈现给全组。

\item
  发起人呈现提案之后，
  在任何讨论之前先进行一次初步表决：
  \begin{itemize}
  \item 谁喜欢这个提案？
  \item 谁能接受这个提案？
  \item 谁对这个提案感到不舒服？
  \end{itemize}
  初步表决可以（线下）用竖大拇指、横挂拇指或倒竖拇指，
  也可以（线上会议）在聊天里输入 +1、0 或 -1。

\item
  如果全组或大多数人喜欢、或能接受这个提案，
  就立即进入正式表决，不再进一步讨论。

\item
  如果大多数人对提案感到不舒服，
  就把它推迟，交回发起人进一步修改。

\item
  如果只有一些成员感到不舒服，他们可以简要陈述自己的反对意见。
  然后设定一个计时器，由引导者主持一场简短的讨论。
  十分钟之后、或没人再有要补充的时候（以先到者为准），
  引导者就下面这个问题请大家做「是/否」表决：
  「我们是否应当在已陈述的反对意见之下实施这个决定？」
  如果多数投「是」，就实施该提案。
  否则，提案退回发起人进一步修改。

\end{enumerate}

\seclbl{线上会议}{s:meetings-online}

\hreffoot{https://chelseatroy.com/2018/03/29/why-do-remote-meetings-suck-so-much/}{Chelsea Troy 关于线上会议为何常常令人沮丧且毫无产出}的讨论，
提出了一个重要观点：
在大多数线上会议里，
在停顿期间第一个开口的人就获得了发言权。
结果呢？
「如果你有想说的话，
你就得停止听正在发言的人，
转而盯着他什么时候会停顿或结束，
好让你能蹿进那一纳秒的寂静、第一个说出点什么。
这种形式{\ldots}鼓励想发言的参与者多说、少听。」

解决办法是在视频会议旁边开一个文字聊天，
让人们可以在那里表示想发言，
然后由主持人从等候名单里挑人。
如果会议规模大、或气氛好争，
就让所有人都静音，只允许主持人给某些人取消静音。

\seclbl{事后复盘}{s:meetings-post-mortem}

每个项目都应当以一次事后复盘结束，
参与者在其中反思刚完成的事，
以及下次可以做得更好的地方。
它的目的\emph{不是}把羞辱的矛头指向个人，
尽管如果真的需要，
事后复盘也是做这件事最合适的地方。

事后复盘像任何其他会议一样进行，
只是多了几条额外的指导原则~\cite{Derb2006}：

\begin{description}

\item[找一位没参与该项目、]
  也与它没有利害关系的主持人。

\item[留出正好一个小时，而且只留一个小时。]
  依我的经验，
  在任何人第一次事后复盘的头十分钟里，
  都听不到什么有用的东西，
  因为人对赞美或痛骂自己的工作都会有点不自在。
  同样，
  过了头一个小时也听不到什么有用的东西：
  如果你还在说，
  那大概是因为一两个人有话想倒出来，
  而不是因为有改进的建议。

\item[要求出席。]
  项目里的每个人都应该在事后复盘现场。
  这比你可能以为的更重要：
  最能从事后复盘中学习的人，
  往往在会议可选时最不可能出现。

\item[列两份清单。]
  我主持时，
  会在白板上写「继续保持」和「下次改变」两个标题，
  然后依次请每个人给每份清单各提一条，
  不得重复已经说过的东西。

\item[评论行为，而不是针对个人。]
  到项目结束时，
  有些人可能已经不是朋友了。
  别让这件事把会议带偏：
  如果某人对团队里另一个人有具体的抱怨，
  就要求他去批评某个具体事件或决定。
  「他态度很差」对任何人改进都\emph{没有}帮助。

\item[给建议排优先级。]
  等每个人的想法都摆到台面上之后，
  按哪些最值得继续保持、
  哪些最需要改变来排序。
  下个项目里你大概只能从每份清单里各着手一两条，
  但如果你每次都这么做，
  你的日子很快就会变得更好。

\end{description}
